\section{Introduction}

The Internet of Things (IoT) is transforming our workflow, and interact with our environment by embedding connectivity into everyday tasks. By enabling these devices to communicate and exchange data seamlessly, IoT is reshaping industries and influencing the future of connectivity. From smart homes and autonomous vehicles to precision agriculture and industrial automation.

By connecting billions of devices globally, IoT creates opportunities to improve efficiency, scalability, and personalization. Its impact extends across industries, fostering interconnected systems where data can support more informed and timely decision-making.

High-mobility scenarios, such as drones and autonomous vehicles, highlight the challenges of maintaining reliable communication in real time. These applications are particularly affected by imparments like the Doppler effect and rapid channel variations, which can affect the stability and accuracy of communication links. Addressing these challenges requires ongoing advancements in wireless communication technologies and network design to meet the demand for low-latency and reliable communication. With its ability to improve productivity, optimize resource utilization, and address specific use cases, IoT is steadily becoming an important element of a more connected and automated future. 

On the other hand, neural networks are following the same trend in the industry with more networks over the market for specifics tasks. However, deploying neural networks in resource-constrained environments, such as embedded systems, presents challenges related to computational complexity, memory usage, and power consumption. Solutions for embedded systems have evolved to address these issues, including the use of optimized hardware accelerators, lightweight neural network architectures, and advanced power management techniques. Tools like TensorFlow Lite and PyTorch Mobile empower developers to deploy neural networks on devices with limited resources, ensuring that IoT solutions can operate effectively at the edge.

With this context, the objective of this research is to conduct experiments aimed at achieving effective classification of channel impairments and cleaning Doppler-shifted signals using neural networks. The proposed approach spans from the physical layer to applications, seeking the best-fit solution through a mixture of experts. This framework will also incorporate constraints related to channel and hardware limitations to ensure practical applicability. The research will initially focus on controlled environments in simulated scenarios. Following this, the solution will be deployed on embedded devices for integration testing, where its performance and validity will be thoroughly evaluated.